\section{Related Work}
\label{sec:related}

\paragraph{Evaluation granularity and localized collateral damage.}
LLM-unlearning benchmarks evaluate several facets of the forget--retain
trade-off, including utility, privacy leakage, and metric robustness
\cite{maini2024tofu,shi2024muse,dorna2025openunlearning}. Methods and
checkpoints, however, are commonly compared through set-level aggregates, and
metric aggregation or model selection can change which method appears
strongest \cite{thaker2024position,dorna2025openunlearning}. Fine-grained
studies show that damage can concentrate in syntactically, semantically, or
representationally adjacent retain regions and can include benign failures
missed by static benchmarks
\cite{chang2025retain,cheng2026entanglement,ko2025holes}. These studies
establish that collateral damage is structured, but mainly localize it after
unlearning or within predefined neighborhoods. Parameter-level localization
and checkpoint differencing are also primarily retrospective: identifying
where knowledge resides or what changed does not by itself predict the outcome
of a future intervention
\cite{hase2023does,chang2024localization,kassem2025mneme}. We instead rank a
declared universe of retained behaviors before the target trajectory.

\paragraph{Prospective auditing.}
Prior trajectory analyses require a realized update to quantify interference
\cite{wang2025gradient}. Closest to our prospective prediction
task, PreUnlearn trains supervised regressors on outcomes from completed runs
to predict a normalized damage ratio for each forget--evaluation-set pair from
pre-update cross-set features \cite{su2026preunlearn}. A related proactive
study predicts algorithm-dependent corruption from model weights and
activations \cite{su2026corruption}. These works establish that
pre-unlearning risk is predictable.

Our distinction is the unit, signal, and downstream test. We rank individual
retained behaviors rather than evaluation sets. The target profiler evaluates
two checkpoint-derived coordinates: a forward-only estimate of Loss
Susceptibility and request exposure in hidden-state space. We do not train
a multivariate target auditor; a single
mixture weight is calibrated on target-request-disjoint development
trajectories and then sealed. Most importantly, we separately test whether the
same rank improves a controlled repair-allocation decision.

\paragraph{Unlearning objectives and retain-side mitigation.}
Approximate objectives suppress output likelihoods or preferences
\cite{jang2023knowledge,yao2024large,zhang2024negative,fan2024simplicity}, or
intervene directly on internal representations, as in RMU
\cite{li2024wmdp}. Representation-control methods have also been adapted as
unlearning baselines \cite{rosati2024repnoise,zou2024circuit}. These labels
describe where an objective reads the model, but do not partition causal
mechanisms: output objectives move hidden states, and representation
objectives ultimately change outputs. We therefore evaluate both profile
coordinates for every parent rather than route parents to a signal endpoint.

Retain-aware methods mitigate damage by reweighting forget examples,
constraining the forget--retain trade-off, curating forget-set coresets, or
regularizing selected retain neighborhoods
\cite{niu2025guard,chen2026hamu,patil2025upcore,chang2025retain,
cheng2026entanglement}. Other methods rectify or project unlearning gradients
against retain-gradient directions
\cite{wang2025gru,huang2024unified,zhou2025implicit,zhou2025geometric}.
Difficulty-aware work similarly uses memorization and entanglement to route
forget examples or adapt the unlearner
\cite{zhao2024what,cheng2026cud,chen2026hamu}. Those decisions alter the
unlearning data or optimization. Our protection experiment leaves the
deletion request and parent trajectory unchanged, begins from a common
criterion-reaching checkpoint, and varies only which retained discovery
behaviors receive a fixed repair quota.

\paragraph{From gradient influence to Loss Susceptibility.}
Influence and scalable data-attribution methods relate training examples to
model behavior through gradients or curvature
\cite{koh2017influence,pruthi2020tracin,park2023trak,kwon2024datainf}, and
gradient features support targeted data selection \cite{xia2024less}.
Forward-only variants reduce candidate-side differentiation, but estimate
training-data value rather than future retained-behavior damage
\cite{ko2024mirrored,deng2026forvalue}. Symmetric perturbations and squared
random projections are classical tools for recovering local gradient
information from function evaluations
\cite{spall1992spsa,nesterov2017random,malladi2023mezo,
hutchinson1989stochastic}. We repurpose this relation to estimate gradient-based
Loss Susceptibility as a pool-shared pre-unlearning diagnostic. Batched forward
evaluations estimate the declared block-local quantity without conditioning on
the direction of the later unlearning update; the resulting coordinate is combined
with request-specific Representation Proximity in one score used for
prediction and protection.

% Keep the overview figure before Section 3. Because this is a two-column
% figure*, the barrier prevents the following section heading from overtaking
% the deferred float.
\begin{figure*}[!t]
\centering
\begingroup
\definecolor{fTwoBlue}{HTML}{0072B2}
\definecolor{fTwoOrange}{HTML}{D55E00}
\definecolor{fTwoGreen}{HTML}{009E73}
\definecolor{fTwoGray}{HTML}{707070}
\definecolor{fTwoLightGray}{HTML}{BDBDBD}
\definecolor{fTwoGrid}{HTML}{D9D9D9}
\definecolor{fTwoInk}{HTML}{1A1A1A}

\begin{tikzpicture}[
  x=1cm,
  y=1cm,
  line cap=round,
  line join=round,
  >={Stealth[length=2.0mm,width=1.3mm]},
  every node/.style={font=\sffamily},
  panelbadge/.style={
    fill=fTwoInk,
    text=white,
    rounded corners=.7mm,
    minimum width=4.2mm,
    minimum height=4.2mm,
    inner sep=0pt,
    font=\sffamily\bfseries\fontsize{7.2}{7.5}\selectfont
  },
  candidate/.style={
    circle,
    draw=fTwoGray,
    fill=white,
    line width=.55pt,
    minimum size=2.35mm,
    inner sep=0pt
  },
  pooloff/.style={
    circle,
    draw=fTwoLightGray,
    fill=white,
    line width=.55pt,
    minimum size=2.05mm,
    inner sep=0pt
  },
  poolon/.style={
    circle,
    draw=none,
    minimum size=2.20mm,
    inner sep=0pt
  }
]
\path[use as bounding box] (0,0) rectangle (17.35,7.65);

% --------------------------------------------------------------------------
% Figure heading
\node[
  anchor=west,
  text=fTwoInk,
  font=\sffamily\bfseries\fontsize{13.2}{13.8}\selectfont
] at (0.20,7.45)
  {Two complementary coordinates for prospective retained-behavior risk};

\node[
  anchor=west,
  text=fTwoGray,
  font=\sffamily\fontsize{7.6}{8.2}\selectfont
] at (0.20,7.02)
  {One frozen two-signal profile, one sealed score, two tests.};

% ==========================================================================
% A. PROFILE
\node[panelbadge] at (0.48,6.43) {A};

\node[
  anchor=west,
  text=fTwoInk,
  font=\sffamily\bfseries\fontsize{9.4}{10.0}\selectfont
] at (0.82,6.43)
  {TWO-SIGNAL PROFILE};

\node[
  anchor=west,
  text=fTwoGray,
  font=\sffamily\fontsize{6.5}{7.1}\selectfont
] at (0.82,6.09)
  {two coordinates, frozen at $\theta_0$};

% Joint-profile field
\shade[
  shading=axis,
  left color=white,
  right color=fTwoGreen!20,
  shading angle=45
] (0.35,1.00) rectangle (7.35,5.72);

\draw[fTwoGrid,dashed,line width=.50pt]
  (3.85,1.00)--(3.85,5.72);
\draw[fTwoGrid,dashed,line width=.50pt]
  (0.35,3.36)--(7.35,3.36);
\draw[fTwoLightGray,line width=.70pt]
  (0.35,1.00) rectangle (7.35,5.72);

% Signal definitions
\node[
  anchor=west,
  text=fTwoBlue,
  font=\sffamily\bfseries\fontsize{7.7}{8.3}\selectfont
] at (0.58,5.39)
  {$\widetilde q_G$: rank-normalized Loss Susceptibility};

\node[
  anchor=west,
  text=fTwoOrange,
  font=\sffamily\bfseries\fontsize{7.7}{8.3}\selectfont
] at (0.58,5.08)
  {$\widetilde q_H$: rank-normalized Representation Proximity};

\node[
  anchor=west,
  fill=white,
  fill opacity=.82,
  text opacity=1,
  inner sep=.45mm,
  text=fTwoInk,
  font=\sffamily\fontsize{7.8}{8.4}\selectfont
] at (0.55,4.66)
  {$S_\alpha=(1-\alpha)\widetilde q_G+\alpha\widetilde q_H$};

% Illustrative retained candidates
\foreach \x/\y in {
  1.33/1.74,
  2.24/2.34,
  3.09/1.84,
  1.55/4.02,
  2.82/3.94,
  4.58/1.92,
  4.66/4.08,
  6.23/3.80,
  6.62/4.18
}{
  \node[candidate] at (\x,\y) {};
}

% Matched Representation Proximity, different Loss Susceptibility
\draw[
  fTwoGray,
  densely dashed,
  line width=.65pt
] (5.40,2.27)--(5.40,4.77);

\node[
  circle,
  draw=fTwoBlue,
  fill=white,
  line width=1.05pt,
  minimum size=3.05mm,
  inner sep=0pt
] at (5.40,2.27) {};

\node[
  circle,
  draw=fTwoBlue,
  fill=fTwoBlue,
  line width=.70pt,
  minimum size=3.05mm,
  inner sep=0pt
] at (5.40,4.77) {};

\node[
  anchor=west,
  text=fTwoBlue,
  font=\sffamily\fontsize{6.7}{7.3}\selectfont
] at (5.65,2.27)
  {$x_{\mathrm{low}}$: lower $\widetilde q_G$};

\node[
  anchor=west,
  text=fTwoBlue,
  font=\sffamily\fontsize{6.7}{7.3}\selectfont
] at (5.65,4.77)
  {$x_{\mathrm{high}}$: higher $\widetilde q_G$};

\node[
  rotate=90,
  anchor=south,
  text=fTwoGray,
  font=\sffamily\fontsize{5.9}{6.5}\selectfont
] at (5.20,3.52)
  {matched Representation Proximity $\widetilde q_H$};

% Increasing two-signal score
\draw[
  -{Stealth[length=2.25mm,width=1.50mm]},
  fTwoGreen,
  line width=1.30pt
] (1.18,1.47)--(6.90,5.25);

\node[
  rotate=34,
  fill=white,
  fill opacity=.88,
  text opacity=1,
  inner sep=.45mm,
  text=fTwoGreen,
  font=\sffamily\bfseries\fontsize{6.8}{7.4}\selectfont
] at (3.83,3.28)
  {two-signal score $S_\alpha$};

\node[
  anchor=north east,
  align=right,
  text=fTwoGreen,
  font=\sffamily\bfseries\fontsize{6.4}{7.0}\selectfont
] at (7.15,5.55)
  {higher $S_\alpha$};

% Region descriptions
\node[
  anchor=south west,
  text=fTwoGray,
  font=\sffamily\fontsize{6.1}{6.7}\selectfont
] at (0.63,1.18)
  {lower $S_\alpha$};

\node[
  anchor=north west,
  text=fTwoGray,
  font=\sffamily\fontsize{6.1}{6.7}\selectfont
] at (0.63,3.28)
  {high $\widetilde q_G$, low $\widetilde q_H$};

\node[
  anchor=south east,
  text=fTwoGray,
  font=\sffamily\fontsize{6.1}{6.7}\selectfont
] at (7.08,1.18)
  {low $\widetilde q_G$, high $\widetilde q_H$};

% Axes
\node[
  anchor=north,
  text=fTwoOrange,
  font=\sffamily\fontsize{7.0}{7.6}\selectfont
] at (3.85,0.72)
  {Representation Proximity $\widetilde q_H$
   \quad low $\longrightarrow$ high};

\node[
  rotate=90,
  anchor=south,
  text=fTwoBlue,
  font=\sffamily\fontsize{7.0}{7.6}\selectfont
] at (-0.02,3.36)
  {Loss Susceptibility $\widetilde q_G$
   \quad low $\longrightarrow$ high};

% ==========================================================================
% B. PREDICTION
\node[panelbadge] at (9.72,6.43) {B};

\node[
  anchor=west,
  text=fTwoInk,
  font=\sffamily\bfseries\fontsize{9.4}{10.0}\selectfont
] at (10.06,6.43)
  {PREDICT: prospective ordering};

\node[
  anchor=west,
  text=fTwoGray,
  font=\sffamily\fontsize{6.5}{7.1}\selectfont
] at (10.06,6.09)
  {$\widehat\alpha_{\mathrm{pred}}$ fixed before target outcomes};

\node[
  anchor=north west,
  align=left,
  text=fTwoInk,
  font=\sffamily\fontsize{6.7}{7.3}\selectfont
] at (9.82,5.79)
  {tests prospective ordering\\and harmful-tail recall};

% Damage outcomes opened only after the prospective rank is sealed
\draw[fTwoInk,line width=.72pt]
  (9.95,4.16)--(17.15,4.16);

\foreach \x/\y in {
  10.12/4.30,
  10.62/4.40,
  11.12/4.35,
  11.62/4.53,
  12.12/4.48,
  12.62/4.66,
  13.12/4.59,
  13.62/4.83,
  14.12/4.75,
  14.62/4.97
}{
  \draw[fTwoGray!72,line width=.68pt]
    (\x,4.16)--(\x,\y);
  \fill[fTwoGray] (\x,\y) circle[radius=1.02mm];
}

\foreach \x/\y in {
  15.12/5.02,
  15.62/4.94,
  16.12/5.27,
  16.72/5.43
}{
  \draw[fTwoOrange!72,line width=.68pt]
    (\x,4.16)--(\x,\y);
  \fill[fTwoOrange] (\x,\y) circle[radius=1.02mm];
}

\draw[fTwoOrange,line width=.90pt]
  (15.02,5.66)--(16.82,5.66);
\draw[fTwoOrange,line width=.90pt]
  (15.02,5.59)--(15.02,5.73);
\draw[fTwoOrange,line width=.90pt]
  (16.82,5.59)--(16.82,5.73);

\node[
  anchor=south,
  text=fTwoOrange,
  font=\sffamily\bfseries\fontsize{6.5}{7.1}\selectfont
] at (15.92,5.71)
  {later damage tail};

\node[
  anchor=south east,
  text=fTwoGray,
  font=\sffamily\fontsize{5.9}{6.5}\selectfont
] at (17.12,4.25)
  {realized damage $d_{t^\dagger}(x)$, observed later};

% Sealed prospective rank
\draw[
  -{Stealth[length=2.0mm,width=1.3mm]},
  fTwoGreen,
  line width=1.05pt
] (9.98,3.98)--(17.08,3.98);

\node[
  anchor=north west,
  text=fTwoGreen,
  font=\sffamily\fontsize{6.6}{7.2}\selectfont
] at (9.98,3.90)
  {sealed $S_{\widehat\alpha_{\mathrm{pred}}}$ audit rank};

% ==========================================================================
% C. PROTECTION
\node[panelbadge] at (9.72,3.31) {C};

\node[
  anchor=west,
  text=fTwoInk,
  font=\sffamily\bfseries\fontsize{9.4}{10.0}\selectfont
] at (10.06,3.31)
  {PROTECT: constrained allocation};

\node[
  anchor=west,
  text=fTwoGray,
  font=\sffamily\fontsize{6.5}{7.1}\selectfont
] at (10.06,2.98)
  {Top-$K_p$ by $S_{\widehat\alpha_{\mathrm{pred}}}$ on
   $\mathcal C_{\mathrm{disc}}$; fixed operator and repair compute budget};

% Selector labels
\node[
  anchor=east,
  text=fTwoInk,
  font=\sffamily\fontsize{6.5}{7.1}\selectfont
] at (11.15,2.42)
  {$S_{\widehat\alpha_{\mathrm{pred}}}$};

\node[
  anchor=east,
  text=fTwoInk,
  font=\sffamily\fontsize{6.5}{7.1}\selectfont
] at (11.15,2.07)
  {$S_0$ (Loss Susceptibility)};

\node[
  anchor=east,
  text=fTwoInk,
  font=\sffamily\fontsize{6.5}{7.1}\selectfont
] at (11.15,1.72)
  {$S_1$ (Proximity)};

\node[
  anchor=east,
  text=fTwoInk,
  font=\sffamily\fontsize{6.5}{7.1}\selectfont
] at (11.15,1.37)
  {random};

% Same candidate universe for every selector
\foreach \yy in {2.42,2.07,1.72,1.37}{
  \foreach \xx in {
    11.43,11.76,12.09,12.42,12.75,
    13.08,13.41,13.74,14.07,14.40
  }{
    \node[pooloff] at (\xx,\yy) {};
  }
}

% Exactly four selected behaviors in every row
\foreach \xx in {13.41,13.74,14.07,14.40}{
  \node[poolon,fill=fTwoGreen] at (\xx,2.42) {};
}

\foreach \xx in {11.76,12.75,13.74,14.40}{
  \node[poolon,fill=fTwoBlue] at (\xx,2.07) {};
}

\foreach \xx in {12.42,13.41,14.07,14.40}{
  \node[poolon,fill=fTwoOrange] at (\xx,1.72) {};
}

\foreach \xx in {11.43,12.42,13.08,14.07}{
  \node[poolon,fill=fTwoGray] at (\xx,1.37) {};
}

% Shared evaluation protocol
\draw[
  -{Stealth[length=1.8mm,width=1.2mm]},
  fTwoInk,
  line width=.75pt
] (14.66,1.90)--(15.00,1.90);

\node[
  draw=fTwoLightGray,
  fill=fTwoGray!5,
  rounded corners=1.2pt,
  minimum width=2.10cm,
  minimum height=1.48cm,
  inner sep=1.5pt
] at (16.16,1.90) {};

\node[
  align=center,
  text=fTwoInk,
  font=\sffamily\bfseries\fontsize{6.1}{6.7}\selectfont
] at (16.16,2.37)
  {EVALUATE\\held-out retain};

\node[
  align=center,
  text=fTwoInk,
  font=\sffamily\fontsize{5.8}{6.5}\selectfont
] at (16.16,1.88)
  {mean $\cdot$ tail CVaR\\native retain; vs. no repair};

\node[
  text=fTwoOrange,
  font=\sffamily\bfseries\fontsize{6.0}{6.6}\selectfont
] at (16.16,1.36)
  {forget + utility};

\node[
  anchor=west,
  text=fTwoGray,
  font=\sffamily\itshape\fontsize{6.0}{6.6}\selectfont
] at (9.70,0.91)
  {the score is sealed before target outcomes; the two tests are separate};

% ==========================================================================
% Sealed-score chain: PROFILE -> PREDICT -> PROTECT

% A PROFILE -> B PREDICT
\draw[
  -{Stealth[length=2.0mm,width=1.3mm]},
  fTwoGreen,
  line width=1.00pt
] (7.58,4.34)
  to[out=10,in=190] (9.43,4.74);

\node[
  fill=white,
  inner sep=.35mm,
  text=fTwoGreen,
  font=\sffamily\fontsize{6.0}{6.6}\selectfont
] at (8.50,4.94)
  {sealed score};

% B PREDICT -> C PROTECT (same sealed score, unlabeled hand-off)
\draw[
  -{Stealth[length=2.0mm,width=1.3mm]},
  fTwoGreen,
  line width=1.00pt
] (9.72,4.14)--(9.72,3.58);

\end{tikzpicture}
\endgroup

\caption{\textbf{One pre-unlearning two-signal profile, two separately tested
uses (schematic).}
Its coordinates are a forward-only, rank-normalized estimate of
gradient-defined Loss Susceptibility and rank-normalized Representation Proximity. A
target-request-disjoint development fold freezes a single two-signal score before the
target trajectory. Prediction compares the sealed audit rank with damage
revealed only after unlearning; the same score independently defines the Top-$K_p$
repair pool on $\mathcal C_{\mathrm{disc}}$. Target audit outcomes
select neither the weight nor the pool. Selector arms share the repair operator, repair compute budget, and
feasibility contract; the unrepaired checkpoint separately measures net repair
benefit. Candidate positions are illustrative;
quantitative evidence is reported in Tables~\ref{tab:pred-value},
\ref{tab:tail-recovery}, \ref{tab:prot-effect},
\ref{tab:prot-contract}, and~\ref{tab:robustness}.}

\Description{A chained schematic from one frozen profile through prediction
to protection. The profile panel places retained behaviors by a forward-only,
rank-normalized estimate of gradient-defined Loss Susceptibility and by
rank-normalized Representation Proximity. Two candidates
with matched Representation Proximity but different Loss Susceptibility
estimates illustrate the
distinct information in the signals. A single frozen score
supplies the audit rank for prospective prediction and independently the Top-$K_p$
discovery pool for protection. Target audit outcomes select neither the
weight nor the pool. Selector arms share
the repair operator, repair compute budget, forgetting constraints, and utility
evaluation.}

\label{fig:profile-two-uses}
\end{figure*}

